\documentclass[10pt,a4paper]{amsart}
\usepackage{amsmath,amssymb,amsthm,graphicx}
% Retained preamble for the Stacks Project chapter morphisms.tex (Morphisms of Schemes).
% The chapter file pinned-original-source.tex carries only \input{preamble} on its first line, so no
% preamble text lives inside the chapter itself.  The Stacks Project uses one single shared preamble for
% every chapter; the snapshot stages it verbatim in sources/stacks-chow/preamble.tex.  The block below is
% copied mechanically, line for line, from that file's lines 14-18 (the Xy-pic setup) and lines 154-217
% (the native theorem declarations together with the mathematical macros).  Nothing was added, removed or
% redefined.  The wrapper class and the amsmath/amssymb/amsthm/graphicx base packages are supplied by the
% assembler, so the class/\externaldocument part of the shared preamble is not repeated here.
% For commutative diagrams we use Xy-pic
\usepackage[all]{xy}

% We use 2cell for 2-commutative diagrams.
\xyoption{2cell}
% Theorem environments.
%
\theoremstyle{plain}
\newtheorem{theorem}[subsection]{Theorem}
\newtheorem{proposition}[subsection]{Proposition}
\newtheorem{lemma}[subsection]{Lemma}

\theoremstyle{definition}
\newtheorem{definition}[subsection]{Definition}
\newtheorem{example}[subsection]{Example}
\newtheorem{exercise}[subsection]{Exercise}
\newtheorem{situation}[subsection]{Situation}

\theoremstyle{remark}
\newtheorem{remark}[subsection]{Remark}
\newtheorem{remarks}[subsection]{Remarks}

\numberwithin{equation}{subsection}

% Macros
%
\def\lim{\mathop{\mathrm{lim}}\nolimits}
\def\colim{\mathop{\mathrm{colim}}\nolimits}
\def\Spec{\mathop{\mathrm{Spec}}}
\def\Hom{\mathop{\mathrm{Hom}}\nolimits}
\def\Ext{\mathop{\mathrm{Ext}}\nolimits}
\def\SheafHom{\mathop{\mathcal{H}\!\mathit{om}}\nolimits}
\def\SheafExt{\mathop{\mathcal{E}\!\mathit{xt}}\nolimits}
\def\Sch{\mathit{Sch}}
\def\Mor{\mathop{\mathrm{Mor}}\nolimits}
\def\Ob{\mathop{\mathrm{Ob}}\nolimits}
\def\Sh{\mathop{\mathit{Sh}}\nolimits}
\def\NL{\mathop{N\!L}\nolimits}
\def\CH{\mathop{\mathrm{CH}}\nolimits}
\def\proetale{{pro\text{-}\acute{e}tale}}
\def\etale{{\acute{e}tale}}
\def\QCoh{\mathit{QCoh}}
\def\Ker{\mathop{\mathrm{Ker}}}
\def\Im{\mathop{\mathrm{Im}}}
\def\Coker{\mathop{\mathrm{Coker}}}
\def\Coim{\mathop{\mathrm{Coim}}}

% Boxtimes
%
\DeclareMathSymbol{\boxtimes}{\mathbin}{AMSa}{"02}

%
% Macros for moduli stacks/spaces
%
\def\QCohstack{\mathcal{QC}\!\mathit{oh}}
\def\Cohstack{\mathcal{C}\!\mathit{oh}}
\def\Spacesstack{\mathcal{S}\!\mathit{paces}}
\def\Quotfunctor{\mathrm{Quot}}
\def\Hilbfunctor{\mathrm{Hilb}}
\def\Curvesstack{\mathcal{C}\!\mathit{urves}}
\def\Polarizedstack{\mathcal{P}\!\mathit{olarized}}
\def\Complexesstack{\mathcal{C}\!\mathit{omplexes}}
% \Pic is the operator that assigns to X its picard group, usage \Pic(X)
% \Picardstack_{X/B} denotes the Picard stack of X over B
% \Picardfunctor_{X/B} denotes the Picard functor of X over B
\def\Pic{\mathop{\mathrm{Pic}}\nolimits}
\def\Picardstack{\mathcal{P}\!\mathit{ic}}
\def\Picardfunctor{\mathrm{Pic}}
\def\Deformationcategory{\mathcal{D}\!\mathit{ef}}

% --- appended 2026-10-12 (Chief): the project own definitions of the reference/slogan/history
% display environments, copied mechanically from the shared Stacks preamble lines 8-12
% (\usepackage{verbatim} plus the three \newenvironment lines, which the chapter preamble keeps
% verbatim from the project preamble). Nothing added, removed or redefined: these environments
% carry no mathematical content and the project itself typesets them as comments. Without these
% three lines a counted statement whose own text uses them fails with Environment slogan undefined.
\usepackage{verbatim}
\newenvironment{reference}{\comment}{\endcomment}
\newenvironment{slogan}{\comment}{\endcomment}
\newenvironment{history}{\comment}{\endcomment}
\usepackage{cleveref}
\crefname{originallemma}{original lemma}{original lemmas}
\Crefname{originallemma}{Original lemma}{Original lemmas}
\makeatletter
\newlabel{morphisms-lemma-alteration-dimension}{{02JX}{source}{}{sourcepointer-1}{}}
\newlabel{morphisms-lemma-alteration-dimension@cref}{{[originallemma][0][]02JX}{[1][0][]source}{}{}{}}
\makeatother
\crefname{originallemma}{original lemma}{original lemmas}
\Crefname{originallemma}{Original lemma}{Original lemmas}
\makeatletter
\newlabel{definition-proper}{{01W1}{source}{}{sourcepointer-2}{}}
\newlabel{definition-proper@cref}{{[originallemma][0][]01W1}{[1][0][]source}{}{}{}}
\makeatother
\crefname{originallemma}{original lemma}{original lemmas}
\Crefname{originallemma}{Original lemma}{Original lemmas}
\makeatletter
\newlabel{schemes-lemma-specialize-points}{{01J7}{source}{}{sourcepointer-3}{}}
\newlabel{schemes-lemma-specialize-points@cref}{{[originallemma][0][]01J7}{[1][0][]source}{}{}{}}
\makeatother
\crefname{originallemma}{original lemma}{original lemmas}
\Crefname{originallemma}{Original lemma}{Original lemmas}
\makeatletter
\newlabel{lemma-dimension-formula}{{02JU}{source}{}{sourcepointer-4}{}}
\newlabel{lemma-dimension-formula@cref}{{[originallemma][0][]02JU}{[1][0][]source}{}{}{}}
\makeatother
\crefname{originallemma}{original lemma}{original lemmas}
\Crefname{originallemma}{Original lemma}{Original lemmas}
\makeatletter
\newlabel{lemma-ubiquity-Jacobson-schemes}{{02J6}{source}{}{sourcepointer-5}{}}
\newlabel{lemma-ubiquity-Jacobson-schemes@cref}{{[originallemma][0][]02J6}{[1][0][]source}{}{}{}}
\makeatother
\crefname{originallemma}{original lemma}{original lemmas}
\Crefname{originallemma}{Original lemma}{Original lemmas}
\makeatletter
\newlabel{lemma-jacobson-finite-type-points}{{01TB}{source}{}{sourcepointer-6}{}}
\newlabel{lemma-jacobson-finite-type-points@cref}{{[originallemma][0][]01TB}{[1][0][]source}{}{}{}}
\makeatother
\crefname{originallemma}{original lemma}{original lemmas}
\Crefname{originallemma}{Original lemma}{Original lemmas}
\makeatletter
\newlabel{lemma-ubiquity-uc}{{02JB}{source}{}{sourcepointer-7}{}}
\newlabel{lemma-ubiquity-uc@cref}{{[originallemma][0][]02JB}{[1][0][]source}{}{}{}}
\makeatother
\crefname{originallemma}{original lemma}{original lemmas}
\Crefname{originallemma}{Original lemma}{Original lemmas}
\makeatletter
\newlabel{topology-lemma-closed-open-map-specialization}{{0066}{source}{}{sourcepointer-8}{}}
\newlabel{topology-lemma-closed-open-map-specialization@cref}{{[originallemma][0][]0066}{[1][0][]source}{}{}{}}
\makeatother
\title{The dimension inequality for a dominant locally finite type morphism}
\author{The Stacks Project Authors}
\date{stacks-project:morphisms.tex:local-snapshot-20261012; GNU FDL 1.2 or later}
\begin{document}
\maketitle
\setcounter{subsection}{1}
\noindent\textit{Editorial note.} All delivered mathematical text is the author's own, copied from the pinned Stacks Project chapter file morphisms.tex; no statement or proof step was rewritten, reordered or removed.\par\medskip
\noindent\textit{Editorial note.} The chapter file pinned-original-source.tex carries only \textbackslash{}input\{preamble\} on its first line, so no preamble text lives inside the chapter. The retained file sources/stacks-morphisms/retained-preamble.tex is a mechanical, line-for-line copy of the project's single shared preamble as staged in the snapshot. Nothing was added, removed or redefined, and the excerpt serves only as preamble, never as a span coordinate.\par\medskip
\noindent\textit{Editorial note.} Both counted spans start at the \textbackslash{}begin line of their environment, so the statement and the author's proof are delivered whole; surrounding prose of the section is deliberately outside the counted spans and is carried as standard prerequisites instead.\par\medskip
\noindent\textit{Editorial note.} The citation fields (source\_id, source\_version, authors, source\_url, license.name, source\_title) are identical to those of every other card this lane has produced from this chapter, so the batch selector treats them as one source.\par\medskip
\noindent\textit{Allowed standard prerequisites (editorial).}\begin{itemize}
\item The dimension formula for the local rings at a point and its image. Lemma \textbackslash{}ref\{lemma-dimension-formula\}, cited at proof lines 13766-13772. Gives the middle inequality.
\item A chain of specialisations lies in the spectrum of the local ring at its final point. Schemes, Lemma \textbackslash{}ref\{schemes-lemma-specialize-points\}, cited at proof lines 13764-13766. Bounds the length of the chain by the local dimension.
\end{itemize}\medskip
\ifdefined\hypertarget\hypertarget{sourcepointer-1}{}\fi
\noindent\textit{Source reference.} Original lemma 02JX: focused-seeds/AGP/sources/stacks-snapshot-20261011/tags.txt, source lines 3352--3352. The lemma's own canonical tag, quoted so that readers can resolve the label inside the delivered text; it points at the counted statement itself and carries no proof content.\par\medskip
\ifdefined\hypertarget\hypertarget{sourcepointer-2}{}\fi
\noindent\textit{Source reference.} Original lemma 01W1: focused-seeds/AGP/sources/stacks-snapshot-20261011/tags.txt, source lines 2466--2466. Navigation pointer to a statement the counted proof refers to but does not contain; nothing load-bearing is delegated to it.\par\medskip
\ifdefined\hypertarget\hypertarget{sourcepointer-3}{}\fi
\noindent\textit{Source reference.} Original lemma 01J7: focused-seeds/AGP/sources/stacks-snapshot-20261011/tags.txt, source lines 2004--2004. Navigation pointer to a statement the counted proof refers to but does not contain; nothing load-bearing is delegated to it.\par\medskip
\ifdefined\hypertarget\hypertarget{sourcepointer-4}{}\fi
\noindent\textit{Source reference.} Original lemma 02JU: focused-seeds/AGP/sources/stacks-snapshot-20261011/tags.txt, source lines 3349--3349. Navigation pointer to a statement the counted proof refers to but does not contain; nothing load-bearing is delegated to it.\par\medskip
\ifdefined\hypertarget\hypertarget{sourcepointer-5}{}\fi
\noindent\textit{Source reference.} Original lemma 02J6: focused-seeds/AGP/sources/stacks-snapshot-20261011/tags.txt, source lines 3325--3325. Navigation pointer to a statement the counted proof refers to but does not contain; nothing load-bearing is delegated to it.\par\medskip
\ifdefined\hypertarget\hypertarget{sourcepointer-6}{}\fi
\noindent\textit{Source reference.} Original lemma 01TB: focused-seeds/AGP/sources/stacks-snapshot-20261011/tags.txt, source lines 2368--2368. Navigation pointer to a statement the counted proof refers to but does not contain; nothing load-bearing is delegated to it.\par\medskip
\ifdefined\hypertarget\hypertarget{sourcepointer-7}{}\fi
\noindent\textit{Source reference.} Original lemma 02JB: focused-seeds/AGP/sources/stacks-snapshot-20261011/tags.txt, source lines 3330--3330. Navigation pointer to a statement the counted proof refers to but does not contain; nothing load-bearing is delegated to it.\par\medskip
\ifdefined\hypertarget\hypertarget{sourcepointer-8}{}\fi
\noindent\textit{Source reference.} Original lemma 0066: focused-seeds/AGP/sources/stacks-snapshot-20261011/tags.txt, source lines 234--234. Navigation pointer to a statement the counted proof refers to but does not contain; nothing load-bearing is delegated to it.\par\medskip
\begin{lemma}
\label{lemma-alteration-dimension}
Let $f : X \to Y$ be a morphism of schemes. Assume that
\begin{enumerate}
\item $Y$ is locally Noetherian,
\item $X$ and $Y$ are integral schemes,
\item $f$ is dominant, and
\item $f$ is locally of finite type.
\end{enumerate}
Then we have
$$
\dim(X) \leq \dim(Y) + \text{trdeg}_{R(Y)} R(X).
$$
If $f$ is closed\footnote{For example if $f$ is proper, see
Definition \ref{definition-proper}.} then equality holds.
\end{lemma}

\begin{proof}
Let $f : X \to Y$ be as in the lemma.
Let $\xi_0 \leadsto \xi_1 \leadsto \ldots \leadsto \xi_e$ be
a sequence of specializations in $X$. Set $x = \xi_e$
and $y = f(x)$. Observe that $e \leq \dim(\mathcal{O}_{X, x})$
as the given specializations occur in the spectrum of
$\mathcal{O}_{X, x}$, see Schemes, Lemma \ref{schemes-lemma-specialize-points}.
By the dimension formula, Lemma \ref{lemma-dimension-formula},
we see that
\begin{align*}
e & \leq \dim(\mathcal{O}_{X, x}) \\
& \leq \dim(\mathcal{O}_{Y, y}) +
\text{trdeg}_{R(Y)} R(X) - \text{trdeg}_{\kappa(y)} \kappa(x) \\
& \leq \dim(\mathcal{O}_{Y, y}) + \text{trdeg}_{R(Y)} R(X)
\end{align*}
Hence we conclude that $e \leq \dim(Y) + \text{trdeg}_{R(Y)} R(X)$ as desired.

\medskip\noindent
Next, assume $f$ is also closed.
Say $\overline{\xi}_0 \leadsto \overline{\xi}_1 \leadsto \ldots
\leadsto \overline{\xi}_d$ is a sequence of specializations in $Y$.
We want to show that $\dim(X) \geq d + r$.
We may assume that $\overline{\xi}_0 = \eta$ is the generic point of $Y$.
The generic fibre $X_\eta$ is a scheme locally of finite type over
$\kappa(\eta) = R(Y)$. It is nonempty as $f$ is dominant. Hence by
Lemma \ref{lemma-ubiquity-Jacobson-schemes} it is a Jacobson scheme.
Thus by Lemma \ref{lemma-jacobson-finite-type-points}
we can find a closed point $\xi_0 \in X_\eta$ and the extension
$\kappa(\eta) \subset \kappa(\xi_0)$ is
a finite extension. Note that
$\mathcal{O}_{X, \xi_0} = \mathcal{O}_{X_\eta, \xi_0}$ because
$\eta$ is the generic point of $Y$. Hence we see that
$\dim(\mathcal{O}_{X, \xi_0}) = r$ by Lemma \ref{lemma-dimension-formula}
applied to the scheme $X_\eta$ over the universally catenary
scheme $\Spec(\kappa(\eta))$ (see Lemma \ref{lemma-ubiquity-uc})
and the point $\xi_0$. This means that we can find
$\xi_{-r} \leadsto \ldots \leadsto \xi_{-1} \leadsto \xi_0$
in $X$. On the other hand, as $f$ is closed specializations
lift along $f$, see
Topology, Lemma \ref{topology-lemma-closed-open-map-specialization}.
Thus, as $\xi_0$ lies
over $\eta = \overline{\xi}_0$ we can
find specializations $\xi_0 \leadsto \xi_1 \leadsto \ldots \leadsto \xi_d$
lying over $\overline{\xi}_0 \leadsto \overline{\xi}_1 \leadsto \ldots
\leadsto \overline{\xi}_d$. In other words we have
$$
\xi_{-r} \leadsto \ldots \leadsto \xi_{-1} \leadsto \xi_0
\leadsto \xi_1 \leadsto \ldots \leadsto \xi_d
$$
which means that $\dim(X) \geq d + r$ as desired.
\end{proof}

\end{document}
